%% kapa: exam=1 round=26 year=2015 period=1 form=A subject=economics subject_ko=경제원론 seq=15 number=55
%% source: Q-Net 감정평가사 기출문제 (cst00309) artlSeq=5138001 file=extracted/1008/5138001/제26회_감정평가사 1차_기출/제26회_감정평가사 1차_기출/1교시_감평_A형.pdf
%% conversion: kordoc 4.24.51 -> tools/gen1.py (auto); review=proofread-sonnet; flags=none
\begin{question}{55}
한 지역에 동질의 휘발유를 판매하는 두 주유소 A, B가 꾸르노(Cournot)경쟁을 하고 있다. 이 지역의 휘발유에 대한 시장수요함수는 $Q=8{,}000-2P$이고, A와 B의 한계비용은 1,000원으로 일정하며, 고정비용은 없다. 이윤극대화를 추구하는 A와 B의 균형판매량은? (단, $P$는 가격, $Q=Q_{A}+Q_{B}$이며, $Q_{A}$, $Q_{B}$는 각각 A와 B의 판매량이다.)
\begin{choices}
  \item $Q_{A}=1{,}500,\quad Q_{B}=1{,}500$
  \item $Q_{A}=1{,}500,\quad Q_{B}=2{,}500$
  \item $Q_{A}=2{,}000,\quad Q_{B}=2{,}000$
  \item $Q_{A}=2{,}500,\quad Q_{B}=2{,}500$
  \item $Q_{A}=3{,}000,\quad Q_{B}=3{,}000$
\end{choices}
\end{question}
\answer{3}
